\chapter{Mack's model diagnostics}

The source for this chapter is Thomas Mack, ``Measuring the Variability of
Chain Ladder Reserve Estimates,'' Casualty Actuarial Society Forum, Spring
1994, Appendices G and H, pages 157--166. The rank and label distributions
used below are extra null models, not consequences of the chain-ladder moment
assumptions.

\section{Development-factor correlation}

\begin{definition}[Adjacent-column rank statistic]\label{def:mack1994-spearman}\lean{VerifiedReserving.rankOn, VerifiedReserving.adjacentFactorRows, VerifiedReserving.factorRankOn, VerifiedReserving.spearmanAdjacent, VerifiedReserving.developmentCorrelationTest}\leanok
For each adjacent pair of development-factor columns, restrict the earlier
column to the accident years present in both columns, rerank both restricted
columns, and compute Spearman's coefficient. Mack's triangle-wide statistic is
the inverse-null-variance weighted average of these coefficients.
\end{definition}

\begin{definition}[Exact permutation null]\label{def:mack1994-permutation}\lean{VerifiedReserving.spearmanRho, VerifiedReserving.spearmanPermutationMean, VerifiedReserving.spearmanPermutationVariance}\leanok
For two complete, tie-free rank permutations, $\rho$ is their centered-rank
correlation. Its exact null fixes one rank vector and draws the other uniformly
from all permutations.
\end{definition}

\begin{theorem}[Equation (G4) and its exact null moments]\label{thm:mack1994-spearman}\lean{VerifiedReserving.spearmanRho_eq_one_sub_six_mul_sum_sq_div, VerifiedReserving.spearmanPermutationMean_eq_zero, VerifiedReserving.spearmanPermutationVariance_eq}\leanok
\uses{def:mack1994-permutation}
For $q\ge2$, the centered-rank correlation is
\[
  1-\frac{6\sum_i(r_i-s_i)^2}{q^3-q}.
\]
Under the uniform-permutation null its mean is zero and its variance is
$1/(q-1)$.
\end{theorem}
\begin{proof}\leanok Reindex finite sums by permutations, center the ranks,
and evaluate the diagonal and off-diagonal second moments by exchangeability.
\end{proof}

\begin{definition}[Nominal aggregate calibration]\label{def:mack1994-development-nominal}\lean{VerifiedReserving.mackSpearmanNullMomentTargets, VerifiedReserving.developmentCorrelationNominalVariance, VerifiedReserving.developmentCorrelationAccepts}\leanok
The inverse-variance aggregate has Mack's variance target only if its component
coefficients are pairwise uncorrelated. The $0.67$ Normal cutoff is recorded as
a definition, not as an exact rejection-probability theorem.
\end{definition}

\section{Calendar-year labels}

\begin{definition}[Median split and diagonal score]\label{def:mack1994-calendar}\lean{VerifiedReserving.FactorBand, VerifiedReserving.factorBand, VerifiedReserving.calendarDiagonalRows, VerifiedReserving.calendarSmallCount, VerifiedReserving.calendarLargeCount, VerifiedReserving.calendarRetainedCount, VerifiedReserving.calendarYearScore, VerifiedReserving.calendarYearTest}\leanok
Each development-factor column is split into small and large ranks around its
median. A calendar diagonal scores the smaller of its two label counts, and
the test statistic sums those scores.
\end{definition}

\begin{definition}[Exact fair-label null]\label{def:mack1994-calendar-null}\lean{VerifiedReserving.fairSignMass, VerifiedReserving.balancedLabelCount, VerifiedReserving.calendarNullMean, VerifiedReserving.calendarNullSecondMoment, VerifiedReserving.calendarNullVariance, VerifiedReserving.mackCalendarNullMean, VerifiedReserving.mackCalendarNullVariance}\leanok
For $n$ retained labels, the exact null assigns binomial mass
$\binom nl/2^n$ to $l$ large labels and scores $\min(l,n-l)$. The finite mean,
second moment, and variance are defined before their closed forms.
\end{definition}

\begin{theorem}[Calendar equations (H1) and (H2)]\label{thm:mack1994-calendar-moments}\lean{VerifiedReserving.sum_fairSignMass, VerifiedReserving.calendarNullMean_eq_mackCalendarNullMean, VerifiedReserving.calendarNullSecondMoment_closed, VerifiedReserving.calendarNullVariance_eq_mackCalendarNullVariance}\leanok
\uses{def:mack1994-calendar-null}
The fair-label masses sum to one. For every $n$, the finite binomial sums equal
Mack's closed mean and variance formulas (H1) and (H2); the second moment also
has an explicit closed form.
\end{theorem}
\begin{proof}\leanok Pair symmetric binomial terms, telescope the signed
half-sum, and evaluate the first two binomial factorial moments.
\end{proof}

\begin{theorem}[Five-label check]\label{thm:mack1994-five}\lean{VerifiedReserving.calendarNullMean_five, VerifiedReserving.calendarNullSecondMoment_five, VerifiedReserving.calendarNullVariance_five, VerifiedReserving.mackCalendarNullMoments_five}\leanok
\uses{thm:mack1994-calendar-moments}
For Mack's five-label example the exact mean, second moment, and variance are
$50/32$, $90/32$, and $95/256$, and the closed formulas reproduce them.
\end{theorem}
\begin{proof}\leanok Normalize the finite binomial sums and evaluate them.
\end{proof}

\begin{definition}[Nominal calendar aggregation]\label{def:mack1994-calendar-nominal}\lean{VerifiedReserving.calendarYearNominalMean, VerifiedReserving.calendarYearNominalVariance, VerifiedReserving.calendarYearAccepts}\leanok
Mack sums the component means and variances and applies a Normal cutoff. The
variance sum and cutoff are nominal definitions because separate diagonal
scores are described as almost uncorrelated rather than proved independent.
\end{definition}
